\chapter{Numerics: Tolerances, Ties, and Degeneracy}
\label{ch:numerics}

Every chapter so far has ended a proof and then quietly deferred a difficulty. This
chapter collects them.

The material here is usually treated as implementation folklore---the constants a code
carries that no paper explains. It is not folklore. Each of the decisions below has a
reason that can be stated, and stating them is worth a chapter because they are where
correct mathematics turns into incorrect software. A graduate course that stops at the
theorem leaves students to rediscover this layer on their own, generally under
deadline.

\section{The two kinds of tolerance}

The single most useful thing to know about tolerances in this subject is that there are
two kinds and they are not equally delicate.
\index{tolerance!two kinds}

\begin{quote}
\textbf{A steering tolerance} decides what the algorithm tries next: which constraint
enters, which index toggles, which set is guessed. A bad value costs an extra iteration
or a fallback. \emph{It can never cost correctness.}

\textbf{A verdict tolerance} decides what the algorithm returns: the certificate test,
the feasibility check, the infeasibility declaration. It is the only thing standing
between a wrong answer and the caller.
\end{quote}

Almost every tolerance in an active-set code is of the first kind, and the anxiety
practitioners feel about tolerances is largely misdirected onto them. The repair rule
of Chapter~\ref{ch:guessing}, the entering-constraint rule of
Chapter~\ref{ch:walking}, the violator sets of Algorithm~\ref{alg:elim}: all steering.
Set them roughly and move on.

The verdict tolerances deserve the attention---and even they are less delicate than they
look, for the reason given in Chapter~\ref{ch:certificate}: the quantity being
thresholded is \emph{bimodal}. A candidate built from a well-conditioned working set
satisfies the KKT test to a few multiples of machine epsilon; one built on a set that
is not has no reason to come close. The two populations sit orders of magnitude apart
and the threshold lies in empty space between them. The same is true of the
infeasibility verdict of Proposition~\ref{prop:farkas}: infeasibility is
\emph{macroscopic}, its size set by the geometry of the constraints rather than by the
arithmetic, so separating it from rounding is easy.

\emph{When a threshold sits in a gap, its exact value does not matter; when it sits in a
continuum, no value is right.} The useful engineering question is therefore not ``what
tolerance?'' but ``is this quantity bimodal, and if not, can I make it so?''

\begin{example}[Making a threshold non-delicate by choice of coordinates]
Chapter~\ref{ch:conditioning} solved the RMT problem in standardised coordinates
$y = D^{1/2}w$ partly for conditioning. There is a second benefit. In those coordinates
both the primal quantity (a standardised weight) and the dual quantity (a reduced
gradient) live on the $O(1)$ correlation scale, with noise floor $\bar\mu \approx 0.5$,
so a default $\varepsilon = 10^{-6}$ is a standard KKT tolerance with no data-dependent
scaling. The change of variables removes the dependence on the raw return scale that a
covariance formulation would carry. Choosing coordinates so that the tolerances become
absolute is a real design move.
\end{example}

\section{Scale invariance}

A tolerance is meaningless without a scale, and the cheapest way to get one is to make
the algorithm invariant instead.

Chapter~\ref{ch:walking} gave the canonical instance: dividing the violation by
$\norm{c_i}$ in the entering rule~\eqref{eq:select} makes the selection invariant under
$(c_i,b_i)\mapsto(\gamma c_i,\gamma b_i)$ (Proposition~\ref{prop:scale}). Without it, a
constraint written in basis points rather than in units is selected first merely for
having been written differently. The proposition claims only that the first
\emph{choice} is unchanged; measured over three decades of $\gamma$ either side of
unity, the observed behaviour is stronger---an identical number of additions and
removals, so the whole trajectory is unchanged, with the returned minimiser agreeing to
about $3\times10^{-16}$.

The general rule: \emph{before choosing a tolerance for a quantity, ask what the
quantity's units are, and whether the algorithm can be rewritten so that it has none}.

\section{Slope floors and window guards}

Two guards recur in path-tracing codes and they are guarding different things.

\paragraph{The slope floor.} In the event rules of Chapter~\ref{ch:parametric} every
candidate step is a ratio $t = (\text{gap})/(\text{slope})$. A slope of exactly zero
means the event cannot occur; a slope of $10^{-17}$ means it occurs at
$t \approx 10^{17}$, which is noise amplified. But a \emph{tiny but genuine} slope still
crosses a bound given a long enough parameter range, so filtering at the user tolerance
would let variables drift out of bounds.

The resolution is to test slopes against $\sqrt{\varepsilon_{\mathrm{mach}}}$---not
against the user tolerance---discarding only slopes at floating-point noise level. And
the exact level is not delicate, for the bimodality reason again: genuine slopes sit
orders of magnitude above $\sqrt{\varepsilon_{\mathrm{mach}}}$ and round-off noise
orders below, so nothing lands near the threshold.

\paragraph{The window guard.} The segment is valid only for parameter values at or below
the current one, because the path is traced monotonically. Candidate ratios above the
current value are spurious crossings outside the segment and must be discarded. This is
not a numerical guard at all; it is the algorithm's contract, and a code that omits it
will occasionally step backwards and then loop.

\paragraph{Relative versus absolute in the ratio test.} As Chapter~\ref{ch:reading}
noted, the breakpoints of a large path crowd together. The tolerance that orders them
must therefore be \emph{relative} to the scale of the parameter. With it the path comes
back strictly monotone to floating-point precision; with a fixed absolute tolerance,
distinct breakpoints eventually merge and the parameter steps backwards. Same lesson as
Section~2: get the units right first.

\section{Ties: anti-cycling in practice}

On degenerate problems---tied expected returns, duplicated columns, several constraints
active at once---multiple events occur at the same parameter value. A naive ``pick any
maximiser'' rule can cycle.
\index{anti-cycling}

The remedy throughout this book is a \textbf{Bland-style rule}: among all events within
tolerance of the best ratio, select the one with the lowest index, and within an index
the lowest event-type. This makes the choice deterministic and resolves degeneracies one
index per iteration. Termination then follows by the standard argument: finitely many
active-set states, a monotone parameter, and the lowest-index rule excluding the only
failure mode, which is revisiting a state without progress.

Three practical notes.

First, \emph{determinism is itself the point}, quite apart from termination. A solver
that makes an arbitrary choice among tied events is not reproducible across platforms,
BLAS versions, or compiler flags, and debugging it is miserable.

Second, the rule must be applied to events \emph{within a tolerance} of the best ratio,
not to exact ties, since exact ties are not what one observes.

Third, an iteration cap is a \emph{numerical safety net}, not the termination
guarantee. A correct trace empirically runs $O(n)$ iterations, so a cap of, say,
$100(n+p+1)$ is never approached; exceeding it signals a floating-point failure such as
a tolerance-induced cycle, which should be raised loudly rather than allowed to hang.
Do not confuse the cap with the theorem.

\section{Cancellation, and rewriting differences as quotients}

Remark~\ref{rem:sign} of Chapter~\ref{ch:freeblock} gave the pattern in its cleanest
form. A Householder reflection needs $v_1 = (d_2)_1 - \alpha$ with
$\alpha = \pm\norm{d_2}$; choosing the sign that makes this a difference of nearly equal
quantities loses digits catastrophically when $d_2$ is close to a positive multiple of
$e_1$. The cure is algebraic:
\[
  v_1 = \frac{(d_2)_1^2 - \nu^2}{(d_2)_1 + \sign((d_2)_1)\nu}
      = \frac{-\sign((d_2)_1)\,\tau}{|(d_2)_1| + \nu},
  \qquad \tau = \norm{(d_2)_{2:}}^2, \ \nu = \norm{d_2},
\]
in which every operation combines like-signed quantities.

Two observations generalise. First, \emph{a subtraction of nearly equal quantities can
usually be rewritten as a quotient} by multiplying by the conjugate---the same device
that fixes the quadratic formula and $\sqrt{x+1}-\sqrt{x}$. Look for it whenever a
formula contains a difference whose operands you expect to be close.

Second, note the interaction with Proposition~\ref{prop:invariance}: the sign convention
is \emph{immaterial to the iterates}, so one is free to choose it. The rewriting is what
makes it immaterial to the \emph{accuracy} as well. Mathematical freedom and numerical
freedom are not the same freedom, and having established the first does not discharge
the second.

\section{When the iterate leaves the feasible set}

Here is a failure mode that is neither a conditioning failure nor a rank failure in the
usual sense, and that earlier versions of real codes aborted on.

On a near-rank-deficient quadratic form the smallest eigenvalues span near-flat
directions of the objective. Accumulated round-off over the hundreds of steps of a large
trace nudges a free variable a hair---typically $10^{-5}$ to $10^{-3}$---outside its box
along one of those directions. The candidate is therefore \emph{optimal to solver
precision but not exactly feasible}, and a strict feasibility check rejects it.

Because the deviation lies in the flat directions, projecting the candidate back onto
the feasible box, and rescaling to restore the budget, changes the objective only at the
level of the round-off itself. Doing so lets the trace complete with points that remain
optimal: across a controlled rank-deficiency sweep every completed point matches a
reference QP solve in objective to about $10^{-8}$, and a well-posed trace is unchanged
to the last digit, the projection being a no-op there.

\paragraph{The repair must be bounded, and by the right quantity.} The two regimes are
told apart by the \emph{conditioning of the free block}, not by the size of the violation
it happens to produce. While that block is numerically full rank its solve is reliable
and any infeasibility is round-off, which is projected away; once the free set outgrows
the rank of the quadratic form the block is numerically singular and its solve is
unreliable, so projecting whatever it returns could silently yield a suboptimal answer.
A guard therefore declines as soon as the free block's reciprocal condition number falls
below about $10^{-12}$, raising an actionable diagnosis instead of a plausible number.

Why key on conditioning rather than on the violation? \emph{Because the violation is the
residual of a singular solve, and its magnitude varies with the BLAS and LAPACK build,
whereas the conditioning is read from the block's symmetric eigenvalues and is identical
on every platform.} Guard on the reproducible quantity.

This gives a precise guarantee, and it is worth stating in the form a user can rely on:
while the free block stays numerically full rank the trace completes and every point is
optimal; when it does not, the algorithm \emph{declines rather than mislead}. Note also
what the guarantee does not cover---it speaks to the linear algebra, not to event ties,
which are the separate combinatorial degeneracy of Section~4.

\section{A short checklist}

For a reader who will implement one of these methods:

\begin{enumerate}
\item Classify every tolerance as steering or verdict. Spend your attention on the
      verdicts.
\item For each verdict tolerance, plot the distribution of the quantity it thresholds
      over a realistic test set. If it is bimodal, you are done. If it is not, change
      coordinates or change the test.
\item Make the algorithm scale-invariant wherever a normalisation can do it, so that
      fewer tolerances need scales at all.
\item Use a lowest-index tie-break, applied within tolerance, and make the solver
      deterministic.
\item Carry an iteration cap, and treat hitting it as a bug report rather than an
      answer.
\item Check the certificate against the original data, never against the construction
      that produced the candidate.
\item Guard on conditioning, which is reproducible, rather than on residual magnitudes,
      which are not.
\item Decline loudly rather than return quietly when the guard fires.
\end{enumerate}

\begin{notes}
The two-tolerance distinction and the bimodality argument are from the Goldfarb--Idnani note \cite{schmelzer2026quadprog};
the slope floor, window guard, relative ratio tolerance, and feasible-set projection
with its conditioning guard are from the \texttt{cvxcla} software paper \cite{cvxcla};
the standardised-coordinate tolerance argument is \cite{schmelzer2026rmt}. Higham
\cite{higham2002} is the reference for cancellation and for accuracy analysis
generally; Golub and Van Loan \cite{golub2013} for the transformations themselves.
\end{notes}

\begin{exercises}

\begin{exercise}
Classify every tolerance appearing in Chapters~\ref{ch:walking}--\ref{ch:parametric} as
steering or verdict, and justify each classification in one sentence.
\end{exercise}

\begin{exercise}
Show that $\sqrt{x+1}-\sqrt{x}$ loses roughly half the available digits for large $x$,
and rewrite it as a quotient. Relate the rewriting to~\eqref{eq:nocancel}.
\end{exercise}

\begin{exercise}
Explain why $\sqrt{\varepsilon_{\mathrm{mach}}}$ rather than
$\varepsilon_{\mathrm{mach}}$ is the right scale for a slope floor. (Hint: consider the
error in a computed slope that is the difference of two $O(1)$ quantities.)
\end{exercise}

\begin{exercise}
Construct a two-variable example in which a fixed absolute tolerance in the ratio test
merges two distinct breakpoints, and show that a relative tolerance separates them.
\end{exercise}

\begin{exercise}
Give an example of a solver decision that is deterministic in exact arithmetic but
platform-dependent in floating point, and propose a rule that removes the dependence.
\end{exercise}

\begin{exercise}
Argue that projecting a slightly infeasible iterate back onto its box changes the
objective by $O(\epsilon^2\lambda_{\min})$ when the deviation $\epsilon$ lies in a flat
direction, and explain why this justifies the repair of Section~6.
\end{exercise}

\begin{exercise}[Implementation]
Instrument your solver from Chapter~\ref{ch:walking} to record, over a few hundred
random instances, the maximum KKT residual of every accepted candidate. Plot the
histogram on a log scale and confirm the bimodality claimed in
Chapter~\ref{ch:certificate}. Then deliberately feed it a rank-deficient working set and
see where the second mode lands.
\end{exercise}

\begin{exercise}[Implementation]
Build a controlled rank-deficiency sweep: generate covariance matrices of decreasing
numerical rank, trace a frontier for each, and record for every point whether the free
block passed the conditioning guard, whether the projection fired, and the objective gap
against a reference QP solve. Reproduce the two regimes of Section~6.
\end{exercise}

\end{exercises}
